% Libro 2 · Mundo 1 · Acertijos de orden y asignación (app: pu-orden)
\mundo{Acertijos de orden y asignación}{pu-orden}{Cinco personas en una fila, seis alrededor de una mesa, tres amigas con
tres carreras y tres barrios. Estos acertijos no piden cuentas sino \emph{orden}: un dibujo bien hecho y una lista de lo
que se va descartando. Son de las preguntas más seguras del examen para quien tiene método.}{%
\item \textbf{1.1} Orden en línea
\item \textbf{1.2} Orden circular
\item \textbf{1.3} Tablas de doble entrada
\item \textbf{1.4} Condiciones: tiene que, podría, no puede}

% ------------------------------------------------------------------
\leccion{pu-or-1}{Orden en línea}{Juego de la fila que se arma con fichas.}

\begin{reto}
Seis corredores llegan a la meta, sin empates. Ana llegó antes que Beto pero después de Carla. Dani llegó justo después
de Beto. Eva llegó justo antes que Carla. Fito no fue ni el primero ni el último. ¿Cuál afirmación es
\emph{necesariamente} cierta?
\opcionesd{Ana llegó tercera.}{Dani llegó de última.}{Fito llegó cuarto.}{Carla llegó primera.}
\end{reto}

\subsection{La idea: casillas y bloques}
El primer paso es siempre el mismo: dibujar una fila de casillas numeradas y decidir qué significa «antes» (más cerca
de la casilla 1). Luego las pistas se clasifican:
\begin{itemize}
  \item \clave{Fijas}: «Óscar es el último». Se ponen de inmediato.
  \item \clave{Bloques}: «Dani justo después de Beto». Se dibujan como una ficha doble \fbox{B\,D} que se mueve junta.
  \item \clave{Relativas}: «Ana antes que Beto». Forman cadenas: $C < A < B$.
  \item \clave{Negativas}: «Fito no es el primero». Tachan casillas.
\end{itemize}

Se empieza por las fijas y los bloques, que dejan menos libertad, y se termina con las negativas, que descartan.

\figura{fila}{Bloques fijos y dos casillas libres: quedan dos filas posibles.}

\begin{regla}
\begin{enumerate}
  \item Dibuja la fila y fija el sentido.
  \item Coloca fijas y bloques; une las relativas en cadenas.
  \item Si hay dos posibilidades, \emph{escríbelas las dos} y sigue con ambas.
  \item Para «necesariamente cierto», la afirmación debe cumplirse en \emph{todas} las posibilidades.
\end{enumerate}
\end{regla}

\subsection{Ejemplos resueltos}
\begin{ejemplo}[Edificio de cinco pisos]
En un edificio de 5 pisos vive una familia por piso: Gómez, Ruiz, Toro, Vélez y Zea. Los Ruiz viven más arriba que los
Toro y más abajo que los Zea. Los Gómez viven en el piso 1. Los Vélez viven justo encima de los Zea. ¿En qué piso viven
los Ruiz?
\solucion Gómez en 1. Cadena: Toro $<$ Ruiz $<$ Zea, y Vélez justo encima de Zea. Quedan los pisos 2 a 5 para Toro,
Ruiz, Zea, Vélez, en ese orden obligado: Toro 2, Ruiz 3, Zea 4, Vélez 5. Los Ruiz, en el 3.
\end{ejemplo}

\begin{ejemplo}[Dos posibilidades]
Cuatro libros en un estante: Álgebra, Biología, Cálculo y Derecho. Álgebra no está en los extremos. Biología está justo a
la derecha de Cálculo. Derecho está a la izquierda de Álgebra. ¿Qué libro puede estar en el extremo derecho?
\solucion Álgebra en 2 o 3. Si Álgebra está en 2, Derecho en 1 y el bloque C\,B en 3-4: \textsf{D A C B}. Si Álgebra
está en 3, Derecho en 1 o 2 y el bloque C\,B necesita dos casillas juntas libres: con D en 1 quedan 2 y 4 (separadas); con
D en 2 quedan 1 y 4 (separadas). No cabe. Única: \textsf{D A C B}. En el extremo derecho está Biología.
\end{ejemplo}

\begin{trampa}
Confundir «antes» con «justo antes». «Ana llegó antes que Beto» permite que haya gente entre ellos; «justo antes» no.
\end{trampa}

\begin{atajo}
En preguntas de «necesariamente», busca una posibilidad que contradiga cada opción: la que no se deja contradecir es
la respuesta.
\end{atajo}

\taller
\piso{Entrena}{filas}
\begin{ejercicios}
\item En una fila de 5, Juan está en el centro, Karen justo delante de Juan, Lina de última y Mario delante de Nelly. ¿Quién es el primero? \sol{Posiciones 1 a 5: Karen 2, Juan 3, Lina 5. Mario y Nelly ocupan 1 y 4, con Mario delante: Mario 1.}
\item En una carrera de 4, Pía no ganó, Quino llegó justo después de Rita y Sara llegó antes que Pía pero después de Quino. ¿Quién ganó? \sol{Cadena R Q … S … P: el orden es R, Q, S, P. Ganó Rita.}
\item Cinco casas en una calle, numeradas de izquierda a derecha. La azul está entre la roja y la verde (no necesariamente juntas). La blanca está en un extremo y la amarilla junto a la blanca. La roja está en la casa 3. ¿Dónde puede estar la verde? \sol{Blanca en 1 o 5 y amarilla a su lado. Si blanca 1, amarilla 2, roja 3: azul y verde en 4 y 5 con la azul entre roja y verde: azul 4, verde 5. Si blanca 5, amarilla 4: azul 2 y verde 1. Verde en 1 o en 5.}
\end{ejercicios}
\piso{Como en el examen}{opción múltiple}
\begin{ejercicios}
\item Con las pistas del ejercicio 3, ¿cuál es necesariamente cierta?
\opcionesd{La azul está junto a la roja.}{La verde está en la casa 5.}{La blanca está en la casa 1.}{La amarilla está en la casa 2.}
\sol{A. En los dos casos la azul (4 o 2) está junto a la roja (3).}
\item Seis estudiantes presentan exposiciones, una por hora. Tomás expone antes que Ulises; Valeria expone justo después de Tomás; Walter expone de tercero. Si Ulises expone de quinto, ¿en qué hora puede exponer Tomás?
\opciones{solo en la 1}{en la 1 o en la 4}{solo en la 2}{en la 1 o en la 2}
\sol{A. Walter ocupa la hora 3 y Ulises la 5. El bloque T\,V necesita dos horas seguidas libres antes de la 5: solo 1-2 (con 4-5 Valeria chocaría con Ulises). Tomás expone en la 1.}
\end{ejercicios}

\begin{retoresuelto}
Bloques \fbox{E\,C} y \fbox{B\,D}, y la cadena $C < A < B$: \textsf{E C … A … B D}. Fito no puede ir en los extremos, así
que entra entre C y A o entre A y B: \textsf{E C F A B D} o \textsf{E C A F B D}. En ambos, Dani es el último.
Ana puede ser 3.ª o 4.ª y Fito 3.º o 4.º; Carla es 2.ª. \textbf{Opción B.}
\end{retoresuelto}
\bitacora

% ------------------------------------------------------------------
\leccion{pu-or-2}{Orden circular}{Juego de la mesa redonda.}

\begin{reto}
Seis amigos se sientan alrededor de una mesa redonda con seis puestos iguales. Ana está frente a Beto. Carlos se sienta
junto a Ana. Elena está frente a Carlos. Diana se sienta junto a Beto. ¿Quién está frente a Diana?
\opciones{Ana}{Elena}{Fabio}{Carlos}
\end{reto}

\subsection{La idea: fija a una persona}
En una mesa redonda, girar a todos un puesto no cambia nada: las relaciones («frente a», «junto a», «a la derecha de»)
siguen iguales. Por eso el primer paso es \clave{fijar} a una persona en un puesto y ubicar a las demás respecto a ella.

\figura{mesa}{Una mesa de seis puestos con una persona fija.}

\subsubsection{Derecha e izquierda}
Si las personas miran hacia el centro de la mesa, la derecha de alguien es el puesto siguiente en el sentido de las
manecillas del reloj… o en el contrario, según cómo esté dibujado. Para no equivocarse: \emph{ponte en su silla}. Mirando
al centro, tu mano derecha queda hacia un lado fijo; en el dibujo de arriba (vista desde el techo), la derecha de quien
está en el puesto 1 (arriba, mirando hacia abajo) es el puesto 6.

\begin{regla}
\begin{itemize}
  \item Fija a la persona con más pistas en el puesto 1.
  \item Con $n$ puestos (par), frente al puesto $k$ está el $k + \tfrac n2$.
  \item Una pista de «junto a» deja dos opciones: si nada más las distingue, por simetría basta estudiar una.
\end{itemize}
\end{regla}

\subsection{Ejemplos resueltos}
\begin{ejemplo}[Cuatro en una mesa cuadrada]
Cuatro amigos, uno en cada lado de una mesa cuadrada: Pedro frente a Quique, Rosa a la derecha de Pedro. ¿Quién está a
la izquierda de Pedro?
\solucion Pedro y Quique enfrentados; Rosa ocupa uno de los otros dos lados (el de la derecha de Pedro). El que queda,
Saúl, está a la izquierda de Pedro.
\end{ejemplo}

\begin{ejemplo}[Ocho puestos]
Ocho personas en una mesa redonda. Entre Juan y Kelly hay exactamente dos personas (por el lado más corto). ¿Cuántas
personas hay entre ellos por el otro lado?
\solucion Además de Juan y Kelly son 6 personas; si 2 están por un lado, 4 están por el otro.
\end{ejemplo}

\begin{trampa}
Fijar a dos personas a la vez en puestos arbitrarios. Solo la primera se puede fijar libremente; las demás se ubican
según las pistas.
\end{trampa}

\taller
\piso{Entrena}{mesas}
\begin{ejercicios}
\item Cinco amigos en una mesa redonda de 5 puestos: Lía junto a Mía y a Nía; Oto junto a Nía. ¿Quién está junto a Pía? \sol{Orden circular Mía, Lía, Nía, Oto, Pía (cerrando con Mía): Pía está junto a Oto y a Mía.}
\item En una mesa de 6, ¿de cuántas maneras puede sentarse una pareja frente a frente si los demás puestos no importan? \sol{Hay 3 pares de puestos enfrentados y la pareja puede intercambiarse: 6 maneras (si los puestos están numerados).}
\end{ejercicios}
\piso{Como en el examen}{opción múltiple}
\begin{ejercicios}
\item Seis personas en una mesa redonda: Q frente a R; S junto a Q y junto a T; U frente a S. ¿Quién está frente a T?
\opciones{Q}{V}{U}{R}
\sol{B. Q en 1, R en 4. S junto a Q: 2 (por simetría). T junto a S: 3. U frente a S: 5. Queda V en 6, frente a T (3).}
\item Ocho caballeros en una mesa redonda; cada uno estrecha la mano solo a quienes \emph{no} tiene al lado. ¿Cuántos apretones hay?
\opciones{20}{28}{16}{24}
\sol{A. Cada uno saluda a 5: $\tfrac{8 \cdot 5}{2} = 20$.}
\end{ejercicios}

\begin{retoresuelto}
Ana en 1, Beto en 4. Carlos junto a Ana: en 2 (el 6 es simétrico). Elena frente a Carlos: 5. Diana junto a Beto: 3 o 5,
y 5 está ocupado: 3. Fabio en 6. Frente a Diana (3) está el 6: Fabio. \textbf{Opción C.}
\end{retoresuelto}
\bitacora

% ------------------------------------------------------------------
\leccion{pu-or-3}{Tablas de doble entrada}{Juego de la tabla con chulos y equis.}

\begin{reto}
Laura, Marta y Nubia estudian Medicina, Derecho e Ingeniería (una carrera cada una) y viven en Bello, Envigado e Itagüí
(un municipio cada una). La de Medicina vive en Envigado. Laura no vive en Bello. Marta vive en Itagüí. La que vive en
Bello no estudia Ingeniería. ¿Qué estudia Nubia?
\opciones{Medicina}{Derecho}{Ingeniería}{No se puede saber}
\end{reto}

\subsection{La idea: una ✓ por fila y por columna}
Cuando hay que emparejar elementos de dos o más grupos, la herramienta es una \clave{tabla de doble entrada}: personas en
las filas, opciones en las columnas. Cada pista pone ✗ o ✓. Como cada persona tiene exactamente una opción y cada
opción una persona, cuando se pone un ✓ se tacha el resto de su fila y de su columna; cuando en una fila queda una sola
casilla libre, esa es el ✓.

\begin{center}\small
\begin{tabular}{l|ccc|ccc}
 & Med & Der & Ing & Bello & Env & Itag\\\hline
Laura & ✓ & ✗ & ✗ & ✗ & ✓ & ✗\\
Marta & ✗ &  &  & ✗ & ✗ & ✓\\
Nubia & ✗ &  &  & ✓ & ✗ & ✗\\
\end{tabular}
\end{center}

Con tres categorías (persona, carrera, municipio), las pistas que mezclan dos categorías («la de Medicina vive en
Envigado») se usan para \emph{transportar} información de una tabla a la otra.

\begin{regla}
\begin{itemize}
  \item Una ✓ por fila y por columna en cada bloque.
  \item Fila o columna con una sola casilla libre: ✓ obligado.
  \item Las pistas negativas son las más útiles: tachan.
  \item Las pistas cruzadas («la de X vive en Y») conectan los bloques.
\end{itemize}
\end{regla}

\subsection{Ejemplos resueltos}
\begin{ejemplo}[Mascotas]
Andrés, Bruno y César tienen un perro, un gato y un loro. Andrés es alérgico a los gatos. El dueño del loro es vecino de
Bruno. César no tiene perro. Bruno no tiene perro. ¿Quién tiene el loro?
\solucion Bruno no tiene el loro (es vecino de su dueño) ni el perro: tiene el gato. César no tiene perro ni gato: tiene
el loro. Andrés, el perro (y no tiene gato, como exige su alergia).
\paso{Comprobación de unicidad.} Sin la última pista habría dos soluciones (César con loro y Andrés con perro, o Andrés con
loro y César con gato). Cuando llegues a dos soluciones, revisa si te faltó usar una pista.
\end{ejemplo}

\begin{trampa}
Usar una pista positiva solo en un sentido. «La de Medicina vive en Envigado» también dice que quien vive en Envigado
estudia Medicina, y que nadie de Bello o Itagüí estudia Medicina.
\end{trampa}

\taller
\piso{Entrena}{tablas}
\begin{ejercicios}
\item Tres hermanos (Iván, José, Kevin) tienen 10, 12 y 15 años. Iván es mayor que José. Kevin no es el mayor. José no es el menor. ¿Cuántos años tiene Kevin? \sol{José no es el menor y es menor que Iván: José 12, Iván 15. Kevin 10.}
\item Cuatro amigas visten de rojo, azul, verde y negro. Ana no viste de rojo ni de negro. Bea viste de verde. Cata no viste de rojo. ¿De qué color viste Dora? \sol{Bea verde. Ana: azul (ni rojo ni negro ni verde). Cata: negro. Dora: rojo.}
\end{ejercicios}
\piso{Como en el examen}{opción múltiple}
\begin{ejercicios}
\item Pablo, Raúl y Saúl son médico, ingeniero y abogado. El ingeniero es el menor de los tres. Pablo es el mayor. Saúl es mayor que Raúl. Saúl no es abogado. ¿Qué es Saúl?
\opciones{médico}{ingeniero}{abogado}{no se puede saber}
\sol{A. Por edades, Pablo $>$ Saúl $>$ Raúl: Raúl es el menor, luego el ingeniero. Saúl no es abogado: es médico (y Pablo abogado).}
\end{ejercicios}

\begin{retoresuelto}
Marta vive en Itagüí; Laura no vive en Bello, así que vive en Envigado; Nubia vive en Bello. Quien vive en Envigado
estudia Medicina: Laura. Quien vive en Bello (Nubia) no estudia Ingeniería, así que estudia Derecho, y Marta Ingeniería.
\textbf{Opción B.}
\end{retoresuelto}
\bitacora

% ------------------------------------------------------------------
\leccion{pu-or-4}{Condiciones: tiene que, podría, no puede}{Juego de los casos que se tachan.}

\begin{reto}
De cinco personas (P, Q, R, S y T) se elige un comité de tres con estas condiciones: si P está, Q no está; R y S están
los dos o ninguno de los dos; T tiene que estar. ¿Cuál comité es posible?
\opciones{P, Q, T}{R, S, T}{P, R, T}{Q, S, T}
\end{reto}

\subsection{La idea: tres preguntas distintas}
En los problemas con condiciones, el examen pregunta casi siempre una de estas cosas:
\begin{itemize}
  \item \clave{Tiene que} (necesariamente): ocurre en \emph{todas} las configuraciones válidas.
  \item \clave{Podría} (es posible): ocurre en \emph{al menos una}.
  \item \clave{No puede}: no ocurre en \emph{ninguna}.
\end{itemize}
Cuando hay pocas configuraciones, lo más seguro es \emph{listarlas todas} y revisar cada opción contra la lista.

\subsubsection{Leer bien el «si…, entonces…»}
«Si P está, Q no está» solo dice algo cuando P está. Si P no está, Q puede estar o no. Tampoco dice «si Q no está,
entonces P está». Lo único equivalente es su \clave{contrarrecíproco}: «si Q está, P no está».

\begin{regla}
\begin{itemize}
  \item «Si A, entonces B» equivale a «si no B, entonces no A». No equivale a «si B, entonces A».
  \item «A o B, pero no ambos»: exactamente uno.
  \item «A solo si B»: si A, entonces B.
  \item Para «podría», basta un ejemplo; para «tiene que», hace falta que no haya contraejemplo.
\end{itemize}
\end{regla}

\subsection{Ejemplos resueltos}
\begin{ejemplo}[Lista de casos]
Con las condiciones del reto, ¿qué persona tiene que estar en el comité, además de T?
\solucion Comités de 3 con T: faltan 2 de P, Q, R, S. Con R y S juntos: \{R, S, T\}. Sin R ni S: dos de P, Q, que no
pueden ir juntos: imposible. Hay un solo comité válido, \{R, S, T\}: R y S tienen que estar.
\end{ejemplo}

\begin{ejemplo}[Horario]
Cuatro materias (A, B, C, D) se dictan de lunes a jueves, una por día. B va antes que D. A no va el lunes. C va el
miércoles o el jueves. ¿Cuál podría ir el martes?
\solucion Listado. Si C es miércoles: lunes, martes y jueves para A, B, D con A no lunes y B antes que D: (B, A, C, D),
(B, D, C, A). Si C es jueves: lunes a miércoles para A, B, D: (B, A, D, C), (B, D, A, C). El martes puede ser A o D; nunca
B ni C.
\end{ejemplo}

\begin{trampa}
Leer «si P, entonces no Q» como «P y Q nunca están juntos, \emph{y} siempre está uno». La condición permite que no esté
ninguno.
\end{trampa}

\taller
\piso{Entrena}{casos}
\begin{ejercicios}
\item Con el horario del ejemplo, ¿qué día va necesariamente B? \sol{El lunes, en los cuatro casos.}
\item En una fiesta, si va Carlos va Diana; si va Diana no va Ernesto. Ernesto fue. ¿Qué se puede concluir? \sol{Diana no fue (contrarrecíproco de la segunda), y entonces Carlos no fue (contrarrecíproco de la primera).}
\end{ejercicios}
\piso{Como en el examen}{opción múltiple}
\begin{ejercicios}
\item Se eligen dos de cuatro libros (W, X, Y, Z). Si se elige W, se elige X. No se eligen Y y Z juntos. ¿Cuál pareja \emph{no} es posible?
\opciones{X, Y}{W, X}{W, Y}{X, Z}
\sol{C. Con W debe ir X, y solo caben dos.}
\item Con las condiciones del ejercicio anterior, ¿cuál afirmación es necesariamente cierta?
\opcionesd{Siempre se elige X.}{W nunca se elige.}{Se elige Y o Z.}{Si no se elige X, se eligen Y y Z.}
\sol{A. Las parejas válidas son \{W, X\}, \{X, Y\} y \{X, Z\}: X está en todas. Sin X, solo quedarían \{Y, Z\} (prohibida) o parejas con W (que exige X).}
\end{ejercicios}

\begin{retoresuelto}
A tiene P y Q juntos; C tiene R sin S; D tiene S sin R. Solo \{R, S, T\} cumple todo. \textbf{Opción B.}
\end{retoresuelto}
\bitacora

% ------------------------------------------------------------------
\evaluacion
\begin{ejercicios}
\item En una fila de 5, Ema está justo detrás de Fer, Gus es el primero, Hilda está detrás de Ema e Iris no es la última. ¿Quién es la última?
\opciones{Hilda}{Ema}{Iris}{Fer}
\sol{A. Gus 1; el bloque F\,E y Hilda detrás de Ema; Iris no última: Gus, Iris, Fer, Ema, Hilda o Gus, Fer, Ema, Iris, Hilda. En los dos, Hilda es última. (1.1)}
\item Cuatro carros en un parqueadero en línea (1 a 4). El rojo está junto al azul. El blanco está en el puesto 1. El negro no está junto al blanco. ¿Dónde está el negro?
\opciones{2}{3}{4}{3 o 4}
\sol{C. Puestos 2, 3, 4 para rojo, azul, negro; negro no en 2. Si negro en 3, rojo y azul en 2 y 4: no son vecinos. Negro en 4. (1.1)}
\item Cinco amigos se sientan en una mesa redonda de 5 puestos. Sofía no está junto a Tomás. ¿Cuántos vecinos tiene Sofía que no sean Tomás?
\opciones{1}{2}{3}{4}
\sol{B. Todos tienen dos vecinos. (1.2)}
\item Seis personas en una mesa redonda: M frente a N; O junto a M; P frente a O; Q junto a N. ¿Quién está frente a Q?
\opciones{M}{R}{O}{P}
\sol{B. M en 1, N en 4, O en 2 (el 6 es simétrico), P en 5. Q junto a N: 3 (el 5 está ocupado). R queda en 6, frente al 3. (1.2)}
\item Tres amigos (X, Y, Z) toman café, té y jugo. X no toma café. Y toma té. ¿Qué toma Z?
\opciones{café}{té}{jugo}{no se sabe}
\sol{A. X toma jugo. (1.3)}
\item Ana, Bea y Cami tienen 17, 18 y 19 años. La de 18 es prima de Ana. Bea es menor que Cami. ¿Cuántos años tiene Ana?
\opciones{17}{18}{19}{no se sabe}
\sol{D. Ana no tiene 18 y Bea es menor que Cami. Si Ana tiene 17: Bea 18 y Cami 19. Si Ana tiene 19: Bea 17 y Cami 18. Las dos cumplen. (1.3)}
\item Si Luis va al paseo, Mario no va. Mario fue al paseo. Entonces
\opcionesd{Luis fue.}{Luis no fue.}{Luis pudo ir o no.}{Nadie más fue.}
\sol{B. Contrarrecíproco. (1.4)}
\item Se escogen 2 de 4 colores (azul, blanco, café, dorado). Si se escoge azul, no se escoge blanco. Café se escoge siempre. ¿Cuántas parejas son posibles?
\opciones{1}{2}{3}{4}
\sol{C. Café con azul, blanco o dorado. (1.4)}
\end{ejercicios}
